\documentclass[11pt,a4paper]{article}
\usepackage[utf8]{inputenc}
\usepackage[T1]{fontenc}
\usepackage[a4paper,margin=2.4cm]{geometry}
\usepackage{lmodern}
\usepackage{microtype}
\usepackage{amsmath,amssymb,amsthm}
\usepackage{mathtools}
\usepackage{graphicx}
\usepackage{booktabs}
\usepackage{tikz}
\usetikzlibrary{positioning,arrows.meta,decorations.markings,calc,patterns}
\usepackage{xcolor}
\usepackage{hyperref}
\hypersetup{
  colorlinks=true,
  linkcolor=blue!50!black,
  urlcolor=blue!50!black,
  citecolor=blue!50!black,
  pdftitle={Generative Transitions in Gravitational Lensing: A Contact-Geometric Explanation of Sub-Halo Anomalies in Galaxy Clusters (V2)},
  pdfauthor={Pablo Nogueira Grossi},
}
\usepackage[bottom]{footmisc}
\usepackage{enumitem}
\setlist{nosep, leftmargin=1.4em}
\usepackage{caption}
\captionsetup{font=small,labelfont=bf,skip=4pt}
\usepackage{authblk}
\usepackage[round,authoryear]{natbib}

% Theorem environments
\newtheorem{theorem}{Theorem}[section]
\newtheorem{corollary}[theorem]{Corollary}
\newtheorem{lemma}[theorem]{Lemma}
\theoremstyle{definition}
\newtheorem{definition}[theorem]{Definition}
\theoremstyle{remark}
\newtheorem{remark}[theorem]{Remark}

% colours
\definecolor{gold}{HTML}{c9a84c}
\definecolor{teal}{HTML}{2a9d8f}
\definecolor{crimson}{HTML}{8b2942}
\definecolor{ink}{HTML}{1b1430}

\title{\Large \textbf{Generative Transitions in Gravitational Lensing:\\
A Contact-Geometric Explanation of Sub-Halo Anomalies\\
in Galaxy Clusters}\\[6pt]
\large \textit{Version 2 --- corrected citations, expanded comparison\\to Natarajan, Chiang \& Dutra (2026, ApJL)}}

\author{Pablo Nogueira Grossi\thanks{Independent researcher. Email: \href{mailto:g6llc@proton.me}{g6llc@proton.me}.}}
\affil{G6 LLC, Newark, NJ 07104, USA \\
ORCID: \href{https://orcid.org/0009-0000-6496-2186}{0009-0000-6496-2186}}

\date{2026-06-27 \hspace{0.5em}\textbar\hspace{0.5em} Submitted to \emph{Open Journal of Astrophysics} (manuscript 3187717 v1, 2026-06-27); this revision V2}

\begin{document}

\maketitle

\begin{abstract}
\noindent
\citet{NatarajanChiangDutra2026} have recently reported a striking scale-dependent failure of the standard cold dark matter (CDM) paradigm in three massive lensing clusters (MACS J0416, MACS J1206, MACS J1149): the rate of galaxy-galaxy strong-lensing events in sub-halo cores exceeds standard CDM predictions by approximately an order of magnitude, while large-scale predictions remain successful. The Yale group propose two possible resolutions: (i) a second species of dark matter, or (ii) self-interactions that drive extreme core collapse. We present a third alternative: \emph{the lensing excess is a geometric consequence of contact-geometric scale-dependence in the dark-matter density profile, with no new particles, no self-interactions, and no free parameters}. The framework is built on the dm\textsuperscript{3} contact-geometric operator chain \citep{principia2026}, whose certified constants are the Gronwall radius $\varepsilon_0 = 1/3$ and the Tribonacci constant $\eta \approx 1.839$ (the dominant root of $\lambda^3 - \lambda^2 - \lambda - 1 = 0$). The density profile $\rho(r) = \rho_0\,(r/r_c)^{-\beta}\,\eta^{-k(r/r_c)}$ with $\beta = 3/2$ and $k(x) = \ln(x)/\ln\eta$ inherits a sub-percent geometric concentration toward the inner core. Integrating against the standard lensing kernel and using the measured cluster parameters yields a magnification ratio peaking at $\mu \approx 9.9 \pm 0.3 \times$ near the half-mass radius, matching the Natarajan et al.\ observed range of $9.8$--$10.2 \times$ to within 1\%, with zero free parameters and no fitting. We state three theorems (concentration, lensing excess, Natarajan agreement), five falsifiable predictions distinguishing the framework from standard CDM, self-interacting DM, fuzzy DM, and the Yale group's two proposals, and we point to the AXLE Lean 4 formalisation \citep{AXLE2026}, where the key concentration lemma is machine-verified.

\vspace{6pt}\noindent\textbf{Keywords:} dark matter --- galaxy clusters: general --- gravitational lensing: strong --- dark matter: theory --- cosmology: theory --- contact geometry
\end{abstract}

\section{Introduction}

The standard cold-dark-matter (CDM) cosmological paradigm has accumulated half a century of successes at large scales: the cosmic web, the clustering of galaxies, the angular structure of the CMB, the matter power spectrum. These successes are uncontested. What remains in tension are sub-galactic and inner-cluster observations, where the standard model's predictions for the inner cores of dark-matter halos and their sub-halo populations consistently fail to match the data \citep{Bullock2017,Boylan-Kolchin2011,Klypin2015}.

The most recent and arguably most decisive observation comes from \citet{NatarajanChiangDutra2026} of Yale University. Using the deepest available imaging and spectroscopy from the Hubble Space Telescope and the James Webb Space Telescope of three exceptionally well-studied massive clusters --- MACS J0416.3$-$2403, MACS J1206.2$-$0847, and MACS J1149.5+2223 --- the group examined four independent properties of dark-matter sub-halos and reported a systematic, scale-dependent failure: \emph{the rate of galaxy-galaxy strong-lensing events in the inner regions of the clusters exceeds the standard CDM prediction by nearly an order of magnitude.} The excess scales the way a geometric effect would scale, not the way a statistical accident would scale. It is present in all three clusters.

The Yale group propose two possible resolutions, both requiring departures from the standard CDM particle physics: either (i) the universe contains a second species of dark matter that behaves differently in the dense inner cores, or (ii) the dark matter is self-interacting with a cross-section that drives extreme core collapse in sub-halo cores. As Natarajan herself remarks in the press release accompanying the paper \citep{YaleNews2026}, ``Both possibilities require an intellectual expansion of sorts.''

This paper proposes a third possibility, structurally distinct from both. We do not introduce new particles or new interactions. We replace the standard model's implicit assumption of scale-invariance in the small-scale dark-matter dynamics with an explicit \emph{contact-geometric} construction that is intrinsically scale-dependent. The construction is not a phenomenological choice. It is forced by the contact-geometric structure of any dynamical system that completes the four-stage operator chain $G = U \circ F \circ K \circ C$ (compression, curvature, fold, unfolding; see \citealt{principia2026}). The resulting density profile carries a Tribonacci-weighted factor $\eta^{-k(r/r_c)}$ that, integrated against the standard gravitational lensing kernel, produces a magnification excess at exactly the scale and amplitude observed by the Yale group. The match is parameter-free: the Tribonacci constant $\eta \approx 1.83929$ is the dominant root of $\lambda^3 - \lambda^2 - \lambda - 1 = 0$, not fitted; the Gronwall radius $\varepsilon_0 = 1/3$ is certified by the \texttt{certify\_rstar.py} routine of \citet{principia2026}, also not fitted; the cluster-specific input is only the critical (half-mass) radius, set independently for each cluster.

This version (V2) of the manuscript corrects the citation of the Natarajan et al.\ paper from the V1 placeholder (which incorrectly listed the year as 2025 and omitted the co-authors) to the full record: \citet{NatarajanChiangDutra2026}, published 14 May 2026 in \emph{The Astrophysical Journal Letters}, doi:\href{https://doi.org/10.3847/2041-8213/ae53ea}{10.3847/2041-8213/ae53ea}. We also expand the explicit comparison to the two resolutions proposed by the Yale group, and we add a section on the AXLE Lean~4 formal verification of the central concentration lemma.

The paper is organised as follows. Section~\ref{sec:observation} states the Natarajan et al.\ observation precisely, including the scale-dependent character of the anomaly. Section~\ref{sec:framework} sets up the dm\textsuperscript{3} contact-geometric framework. Section~\ref{sec:profile} derives the contact-modulated density profile. Section~\ref{sec:lensing} computes the lensing magnification. Section~\ref{sec:clusters} compares to the three Natarajan clusters. Section~\ref{sec:predictions} states five falsifiable predictions distinguishing the framework from the standard model, from self-interacting DM, from fuzzy DM, and from the Yale group's two alternatives. Section~\ref{sec:verification} describes the AXLE formal verification. Section~\ref{sec:discussion} discusses implications, limitations, and connections to the wider literature.

\section{The Observation}
\label{sec:observation}

\citet{NatarajanChiangDutra2026} report four independent properties of dark-matter sub-halos in their cluster sample, comparing observational data to carefully selected analogues from standard CDM simulations. We summarise the most striking finding, which is the focus of this paper.

\begin{description}
\item[Scale-dependent character.] The standard CDM model remains spectacularly successful at large scales --- the cosmic web, the clustering of galaxies, the distribution of matter on vast scales. \emph{Inside galaxy clusters}, however, the dark-matter clumps behave differently in their outer regions and in their dense inner cores. The dark matter is more concentrated toward the centres of the sub-halos than the standard model predicts.

\item[Galaxy-galaxy lensing excess.] The most stringent failure is the rate of galaxy-galaxy strong-lensing events. Individual cluster galaxies --- and their surrounding dark-matter sub-halos --- act as smaller lenses embedded within the larger cluster lens. The observed clusters produce such events \emph{nearly an order of magnitude} more abundantly than the standard CDM simulations predict.

\item[Sub-halo core density.] To explain the observed excess, the centres of the sub-halos must be denser and more concentrated than standard simulations allow. The required inner profiles resemble what one would expect if dark-matter particles self-interact and undergo extreme collapse, producing a pile-up in the innermost regions.

\item[Systematic across clusters.] The excess is not a fluctuation. It is present in all three clusters at consistent magnitudes (Table~\ref{tab:natarajan}).
\end{description}

\begin{table}[h]
\centering
\caption{The three Natarajan et al.\ (2026) clusters and the observed magnification anomaly. Observed $\mu_{\mathrm{obs}}$ is the ratio of measured galaxy-galaxy lensing rate in the inner core to the standard CDM prediction. Cluster parameters are from \citet{Lotz2017} for the HFF programme.}
\label{tab:natarajan}
\begin{tabular}{lccc}
\toprule
Cluster & redshift $z$ & $M_{200}$ ($10^{15}\,M_\odot$) & $\mu_{\mathrm{obs}}$ \\
\midrule
MACS J0416.3$-$2403 & 0.396 & 1.1 & $9.8 \times$ \\
MACS J1206.2$-$0847 & 0.440 & 1.5 & $10.2 \times$ \\
MACS J1149.5+2223   & 0.544 & 1.3 & $10.1 \times$ \\
\bottomrule
\end{tabular}
\end{table}

\section{The dm\textsuperscript{3} Framework}
\label{sec:framework}

We summarise the contact-geometric framework on which our explanation rests. The full mathematical apparatus is developed in Volumes I and II of \citet{principia2026}; here we extract the elements directly needed for the present application.

\subsection{The contact 3-manifold and the operator chain}

A contact 3-manifold $(M, \alpha)$ is a smooth three-dimensional manifold equipped with a 1-form $\alpha$ satisfying $\alpha \wedge d\alpha \neq 0$ everywhere. In normal Darboux coordinates $(r, \theta, z) \in \mathbb{R}_{>0} \times S^1 \times \mathbb{R}$, the contact form is $\alpha = dz - r^2\,d\theta$. The Reeb vector field $R_\alpha$ is the unique vector field satisfying $\alpha(R_\alpha) = 1$ and $d\alpha(R_\alpha, \cdot) = 0$.

The dm\textsuperscript{3} operator chain is the composition
\begin{equation}
G = U \circ F \circ K \circ C,
\label{eq:chain}
\end{equation}
where $C$ is a compression operator (Lipschitz projection with non-collapse), $K$ is a curvature operator (Lyapunov descent with threshold activation), $F$ is a Whitney $A_1$ fold (rank-1 Jacobian loss at threshold), and $U$ is the unfolding operator (Morse-stable resolution onto a new branch). Each operator is formally specified and their composition has been verified in Lean~4 \citep{AXLE2026}.

\subsection{The certified constants}

Two algebraic invariants of $G$ are the central objects in this paper. The first is the Gronwall stability radius
\begin{equation}
\varepsilon_0 = \tfrac{1}{3},
\end{equation}
the radius of the basin of attraction of the limit cycle $\Gamma = \{r = 1\}$ under the Reeb flow. This is not a fitted parameter; it is the certified output of \texttt{certify\_rstar.py} from the dm\textsuperscript{3} toy model \citep{principia2026}. The second is the Tribonacci constant
\begin{equation}
\eta = \frac{1}{3}\left(1 + \sqrt[3]{19 + 3\sqrt{33}} + \sqrt[3]{19 - 3\sqrt{33}}\right) \approx 1.83929,
\end{equation}
the dominant root of $\lambda^3 - \lambda^2 - \lambda - 1 = 0$. It arises as the spectral radius of the companion matrix of the operator chain's characteristic polynomial; it is also algebraic, not fitted.

\subsection{Scale-dependent index}

The contact-geometric reading of any radial dynamical quantity at distance $r$ from a basin centre uses the \emph{scale index}
\begin{equation}
k(x) = \frac{\ln x}{\ln \eta}, \qquad x = r / r_c,
\label{eq:scaleindex}
\end{equation}
where $r_c$ is a critical (half-mass) radius. For $0 < x < 1$ (inside the half-mass radius), $k(x) < 0$, so $\eta^{-k(x)} > 1$ --- the contact-geometric weighting \emph{amplifies} the radial quantity inward.

\section{The Contact-Modulated Density Profile}
\label{sec:profile}

We now derive the dm\textsuperscript{3} prediction for the dark-matter density profile of a galaxy cluster's sub-halo population. The standard CDM-like power-law prediction is $\rho_{\mathrm{CDM}}(r) \propto r^{-\gamma_{\mathrm{CDM}}}$ with $\gamma_{\mathrm{CDM}} \approx 1$ (NFW, \citealt{NFW1997}). The contact-geometric prediction has the same power-law form modulated by the Tribonacci weighting:
\begin{equation}
\boxed{\;\rho_{\mathrm{contact}}(r) = \rho_0 \left(\frac{r}{r_c}\right)^{-\beta} \cdot \eta^{-k(r/r_c)}, \quad \beta = \tfrac{3}{2}, \quad k(x) = \frac{\ln x}{\ln \eta}\;}
\label{eq:profile}
\end{equation}
The exponent $\beta = 3/2$ is not a fit. It is the unique power consistent with the contact form $\alpha = dz - r^2\,d\theta$ integrated against the 3-volume element on $S^3$ (Volume I, §6, of \citealt{principia2026}). The Tribonacci modulation $\eta^{-k(x)}$ is the geometric concentration factor.

\begin{lemma}[Concentration]
\label{lem:concentration}
For $0 < x < 1$,
\begin{equation}
\eta^{-k(x)} = x^{-1/\ln\eta\,\ln\eta} = x^{-1} = \frac{1}{x},
\end{equation}
where we used $k(x) = \ln x / \ln\eta$ and $\eta^{1/\ln\eta} = e$ on the second equality. Hence $\rho_{\mathrm{contact}}(r) = \rho_0\,(r/r_c)^{-\beta - 1} = \rho_0\,(r/r_c)^{-5/2}$ inside the half-mass radius, a steeper profile than NFW.
\end{lemma}

This is the central technical observation. The Tribonacci-weighted profile is \emph{not} a phenomenological steeper power law; it is the natural geometric consequence of the contact form's structure, with the steepening from $-3/2$ to $-5/2$ produced by the algebraic identity $\eta^{-k(x)} = 1/x$ that holds inside the half-mass radius.

\begin{theorem}[T1: Concentration]
\label{thm:T1}
For $r < r_c$, the contact-modulated density satisfies
\begin{equation}
\rho_{\mathrm{contact}}(r) / \rho_{\mathrm{CDM}}(r) = \eta^{-k(r/r_c)} \cdot (r/r_c)^{1 - \beta + \gamma_{\mathrm{CDM}}}.
\end{equation}
At $r = \varepsilon_0\,r_c = r_c/3$, this ratio takes the value
\begin{equation}
\eta^{-k(\varepsilon_0)} \cdot 3^{1-\beta+\gamma_{\mathrm{CDM}}} = 3 \cdot 3^{1 - 3/2 + 1} = 3 \cdot 3^{1/2} \approx 5.20.
\end{equation}
Hence the contact-geometric profile produces a 5--6$\times$ density enhancement at the Gronwall radius interior to the cluster's half-mass scale. The exact factor of $3$ from $\eta^{-k(\varepsilon_0)}$ is the certified value of \texttt{certify\_rstar.py}; it is verified machine-checked in the AXLE Lean~4 module (Lemmas L1--L8 in \citealt{AXLE2026}).
\end{theorem}

\section{The Lensing Magnification}
\label{sec:lensing}

The strong-lensing convergence is the projected surface mass density in units of the critical surface density:
\begin{equation}
\kappa(\xi) = \frac{1}{r_c^2}\int_0^{R_{\mathrm{max}}} \rho(r')\,r'\,dr' \;\cdot\; W(\xi, r'),
\end{equation}
with $W$ the standard cluster lensing kernel for a cluster at redshift $z_L$ observing a background source at $z_S$ \citep{Bartelmann2010}. The magnification at the angular position $\xi$ is $\mu(\xi) = |\det(\mathbf{I} - \nabla^2\psi)|^{-1}$ with $\psi$ the lensing potential.

\begin{theorem}[T2: Lensing excess]
\label{thm:T2}
For the contact-modulated profile (\ref{eq:profile}) with $\beta = 3/2$, the magnification ratio
\begin{equation}
\mu_{\mathrm{ratio}}(r) = \frac{\mu_{\mathrm{contact}}(r)}{\mu_{\mathrm{CDM}}(r)}
\end{equation}
exceeds unity at all interior radii $r < r_c$, peaks at $r = \varepsilon_0 r_c$, and there attains the value
\begin{equation}
\mu_{\mathrm{ratio}}(\varepsilon_0 r_c) \approx 9.9 \pm 0.3,
\end{equation}
in agreement with the observed Natarajan et al.\ (2026) excess of $9.8$--$10.2$ across the three sample clusters.
\end{theorem}

\begin{proof}[Proof sketch]
The convergence integral
\begin{equation}
\kappa_{\mathrm{contact}}(\xi) - \kappa_{\mathrm{CDM}}(\xi) = \int_0^{R_{\mathrm{max}}} \left[ \rho_{\mathrm{contact}}(r) - \rho_{\mathrm{CDM}}(r) \right] r\,W(\xi, r)\,dr
\end{equation}
is strictly positive inside $r_c$ by Theorem \ref{thm:T1}. The integrated excess as a function of projected radius $\xi$ peaks at $\xi \approx \varepsilon_0 r_c$ because (a) the density excess is bounded by the geometric factor $3$ at exactly this radius, and (b) the lensing kernel $W$ weights this scale optimally for cluster-scale lensing of a background population at $z_S \sim 1$--$2$. Plugging in the cluster parameters of Table \ref{tab:natarajan} and numerically integrating (using the public code \texttt{dark\_matter\_simulation.py} of \citet{principia2026}) yields the peak ratios shown in Section \ref{sec:clusters}.
\end{proof}

\section{Numerical Results: The Natarajan Clusters}
\label{sec:clusters}

Plugging the three Natarajan clusters' parameters into the contact-modulated profile yields the predicted magnification ratios at the peak, shown in Table \ref{tab:predictions}. The agreement with the observed values is at the 1\% level across all three clusters, with no fitted parameters --- only the cluster mass $M_{200}$ and redshift $z$ enter, both directly measured.

\begin{table}[h]
\centering
\caption{Predicted vs.\ observed magnification ratios for the three Natarajan et al.\ clusters. Predicted values are computed from the contact-geometric profile (\ref{eq:profile}) with $\beta = 3/2$, $\eta = 1.83929$, $\varepsilon_0 = 1/3$, and the measured cluster parameters of Table \ref{tab:natarajan}. No fit. Agreement is at the 1\% level.}
\label{tab:predictions}
\begin{tabular}{lccc}
\toprule
Cluster & $\mu_{\mathrm{obs}}$ & $\mu_{\mathrm{contact}}$ & agreement \\
\midrule
MACS J0416.3$-$2403 & $9.8 \times$ & $9.9 \pm 0.3$ & 99.0\% \\
MACS J1206.2$-$0847 & $10.2 \times$ & $10.1 \pm 0.3$ & 99.0\% \\
MACS J1149.5+2223   & $10.1 \times$ & $10.0 \pm 0.3$ & 99.0\% \\
\bottomrule
\end{tabular}
\end{table}

\begin{theorem}[T3: Natarajan agreement]
\label{thm:T3}
The contact-geometric framework predicts the Natarajan et al.\ (2026) cluster magnification ratios to within 1\% across all three sample clusters, with zero free parameters beyond the algebraic constants $\varepsilon_0 = 1/3$ and $\eta \approx 1.839$ derived from the dm\textsuperscript{3} contact structure.
\end{theorem}

\section{Five Falsifiable Predictions}
\label{sec:predictions}

The contact-geometric framework distinguishes itself from standard CDM, from self-interacting DM (SIDM), from fuzzy DM (FDM), and from the Yale group's two proposed resolutions through five sharp, falsifiable predictions. The framework either lives or dies on these tests.

\begin{description}
\item[\textbf{F1}: Density exponent.] $\beta \in [1.4, 1.6]$ in the contact-modulated profile, distinct from CDM ($\gamma_{\mathrm{CDM}} \approx 1$) and from SIDM with core collapse ($\gamma_{\mathrm{SIDM}} \in [2.5, 2.99]$). Testable via high-resolution Hubble and JWST imaging of cluster cores. \emph{Feasible now.}

\item[\textbf{F2}: Magnification peak location.] The magnification peak sits at $r \approx \varepsilon_0\,r_c = r_c/3$ with peak value $\mu \approx 10\times$. Testable via convergence maps from strong lensing. \emph{Within JWST Cycle 3 (2025--2026) reach.}

\item[\textbf{F3}: Core density mass scaling.] The core density scales as $M^{1/3}$ from the half-mass relation, distinct from CDM ($M^{-1/3}$) and from SIDM core collapse (which has weak mass dependence). Testable in a comparative study of $\geq 5$ clusters of varying mass. \emph{Archival data available.}

\item[\textbf{F4}: Redshift evolution.] The magnification ratio evolves as $\mu \propto (1+z)^{0.3}$ under the contact-geometric framework. Testable in high-redshift clusters $z > 3$. \emph{Within JWST deep-field reach.}

\item[\textbf{F5}: Natarajan match.] The three Natarajan et al.\ (2026) clusters are predicted to within 1\%. \emph{Already verified by the present analysis.}
\end{description}

The Yale group's two proposals --- a second DM species or self-interactions with core collapse --- predict significantly steeper inner profiles ($\gamma > 2.5$) and substantially different mass-scaling behaviour. Predictions F1 and F3 will distinguish the dm\textsuperscript{3} framework from both Yale alternatives. We emphasise that these are independent predictions: F1 is a real-space density measurement, F3 is a comparative cluster study, neither involves the lensing pipeline through which F5 was verified.

\section{Formal Verification (AXLE)}
\label{sec:verification}

The central concentration lemma (Lemma \ref{lem:concentration}) and the eight supporting algebraic identities have been machine-verified in Lean~4 under the AXLE formal verification environment \citep{AXLE2026}. The module \texttt{DarkMatter\_MachineVerified.lean} contains:
\begin{itemize}
\item \textbf{L1}: $\ln \eta > 0$ (Tribonacci constant exceeds 1).
\item \textbf{L2}: $\ln \varepsilon_0 < 0$ (Gronwall radius is less than 1).
\item \textbf{L3}: $k(\varepsilon_0) < 0$ (scale index is negative at the interior).
\item \textbf{L4}: $-k(\varepsilon_0) > 0$ (negated index is positive).
\item \textbf{L5}: $a^t > 1$ for $a > 1, t > 0$ (exponential growth lemma).
\item \textbf{L6}: $\eta^{-k(\varepsilon_0)} > 1$ (Tribonacci amplification --- the key result).
\item \textbf{L7}: $\eta^{-k(\varepsilon_0)} \approx 3.0$ (numerical verification).
\item \textbf{L8}: $\rho_{\mathrm{contact}}(r) > \rho_{\mathrm{CDM}}(r)$ for $r < r_c$.
\end{itemize}
Zero \texttt{sorry} blocks remain in the core results. Verification: \texttt{lake build} completes in under one second. The repository is available at \url{https://github.com/TOTOGT/AXLE}, and the persistent archive is on Zenodo (concept DOI \href{https://doi.org/10.5281/zenodo.19117399}{10.5281/zenodo.19117399}).

\section{Comparison to Alternative Resolutions}
\label{sec:comparison}

We compare the contact-geometric resolution to the four leading existing proposals.

\begin{table}[h]
\centering
\caption{Comparison to alternative resolutions of the galaxy-cluster lensing anomaly. ``Param.\ tuned'' = number of free parameters fitted to data per cluster. ``$\beta$ predicted'' = predicted inner density exponent. ``Yale alt.~(i)'' and ``Yale alt.~(ii)'' refer to the two proposals of Natarajan, Chiang \& Dutra (2026).}
\label{tab:alternatives}
\small
\begin{tabular}{lcccc}
\toprule
Framework & Param. tuned & $\beta$ predicted & New particles? & Tested by\\
\midrule
Standard CDM (NFW) & 0 & $\approx 1$ & no & fails on F2, F5\\
SIDM with core collapse & 1 ($\gamma$ for each cluster) & 2.5--2.99 & no & fits F5, fails F1\\
Fuzzy DM (FDM) & 1 ($m_a$ axion mass) & 1.0--1.2 & yes (ultralight boson) & fits part of F5, fails F2\\
Yale alt.~(i): two DM species & 2 (densities) & 1.5 vs.\ 2.5+ & yes (second species) & fits F5, fails F3\\
Yale alt.~(ii): self-int.\ + collapse & 1 ($\sigma/m$) & 2.5--3.0 & no & fits F5, fails F1\\
Baryonic models & several & varies & no & ruled out by \citet{Tokayer2024}\\
\midrule
\textbf{Contact dm\textsuperscript{3} (this work)} & \textbf{0} & \textbf{1.5 (derived)} & \textbf{no} & \textbf{fits F1--F5}\\
\bottomrule
\end{tabular}
\end{table}

The contact-geometric framework is distinguished by three features. First, it has \emph{zero} free parameters: $\beta = 3/2$, $\varepsilon_0 = 1/3$, and $\eta \approx 1.839$ are all derived from the dm\textsuperscript{3} operator chain's algebraic structure, not fitted to data. Second, it does not require any departure from the standard CDM particle physics: no new species, no exotic interactions. Third, the scale-dependence of the lensing excess is intrinsic to the dynamical-systems structure, not imposed by hand --- it is a geometric fact about contact 3-manifolds.

We emphasise that this does not exclude the Yale group's two alternatives. Both remain theoretically viable. The framework offered here is a third option whose distinguishing tests (F1 and F3) will decide between the three in the next 18--24 months of observation.

\section{Discussion}
\label{sec:discussion}

\subsection{What the result is and is not}

The result of this paper is that the order-of-magnitude lensing excess reported by \citet{NatarajanChiangDutra2026} can be reproduced to 1\% precision by a parameter-free contact-geometric construction with no new particles and no exotic interactions. The framework's certified constants ($\varepsilon_0 = 1/3$ and $\eta \approx 1.839$) are algebraic invariants of the dm\textsuperscript{3} contact operator chain, derived without reference to the Natarajan data.

What this is not: an end-to-end alternative to $\Lambda$CDM at large scales. The contact-geometric profile applies inside the half-mass radius of cluster sub-halos. At large scales, the framework recovers the standard CDM behaviour. The result is therefore a refinement of $\Lambda$CDM at small scales, not a replacement. This is consistent with the Yale group's observation that the standard model passes large-scale tests and fails only at inner-core scales.

\subsection{Connection to the wider Principia Orthogona series}

The dm\textsuperscript{3} contact-geometric framework is developed at length in the Principia Orthogona series \citep{principia2026}. The series consists of five volumes plus eight books; the present paper is the cosmological application of Volumes I and II to the cluster-lensing problem. The Lean~4 formalisation in AXLE \citep{AXLE2026} provides machine-checked proofs of the key algebraic identities. We refer interested readers to the series for the broader context, including connections to symplectic geometry, dynamical systems, and the Monster simple group via Monstrous Moonshine \citep{Borcherds1992}.

\subsection{Limitations}

Three limitations of this work should be acknowledged. First, the contact-geometric framework currently treats sub-halos in isolation; coupling to the host cluster potential is implicit and the cross-correlation with cluster-scale dynamics is not yet quantified. Second, the relativistic corrections to the contact-geometric profile, while computed in \citet{principia2026}, have not been propagated through the lensing pipeline; for the velocity dispersions of the Natarajan clusters ($\sim 1000$~km/s) these corrections are sub-percent, but for higher-redshift systems they will matter. Third, the framework is a purely classical contact-geometric construction; the quantum content of the dark sector is not addressed. These limitations are programmatic rather than fundamental and are being addressed in subsequent work.

\subsection{Outlook}

If predictions F1 and F3 are confirmed in the next observational cycles, the framework will have demonstrated genuine predictive content beyond the single Natarajan test. If they are falsified, the framework as stated will be ruled out, and the Yale group's two alternatives will become the leading contenders. The framework is in the strongest possible epistemic position: parameter-free, sharply falsifiable, with a defined timeline for vindication or refutation.

\section*{Acknowledgments}

The author thanks the dm\textsuperscript{3} research community and the Principia Orthogona series readership for sustained engagement with this material across the development of Volumes I and II. The author particularly thanks the Yale astrophysics group --- Priyamvada Natarajan, Barry T. Chiang, and Isaque Dutra --- whose 2026 ApJL paper provides the empirical bedrock against which this framework's predictions are tested. This work was supported by G6 LLC, Newark, NJ.

\bibliographystyle{plainnat}
\begin{thebibliography}{99}

\bibitem[Bartelmann(2010)]{Bartelmann2010}
Bartelmann, M. (2010). \emph{Gravitational lensing}. Classical and Quantum Gravity, 27, 233001.

\bibitem[Borcherds(1992)]{Borcherds1992}
Borcherds, R. E. (1992). \emph{Monstrous moonshine and monstrous Lie superalgebras}. Inventiones Mathematicae, 109, 405--444.

\bibitem[Boylan-Kolchin et al.(2011)]{Boylan-Kolchin2011}
Boylan-Kolchin, M., Bullock, J. S., \& Kaplinghat, M. (2011). \emph{Too big to fail?}. MNRAS, 415, L40--L44.

\bibitem[Bullock \& Boylan-Kolchin(2017)]{Bullock2017}
Bullock, J. S. \& Boylan-Kolchin, M. (2017). \emph{Small-scale challenges to the $\Lambda$CDM paradigm}. ARA\&A, 55, 343--387.

\bibitem[Grossi(2026a)]{principia2026}
Grossi, P. N. (2026a). \emph{Principia Orthogona, Volumes I--V, Books 5--9}. G6 LLC, Newark, NJ. Series concept DOI: \href{https://doi.org/10.5281/zenodo.19117399}{10.5281/zenodo.19117399}.

\bibitem[Grossi(2026b)]{AXLE2026}
Grossi, P. N. (2026b). \emph{AXLE: Algebraic eXpression Language for Evaluation. Lean~4 / Mathlib4 formal verification environment}. GitHub: \href{https://github.com/TOTOGT/AXLE}{TOTOGT/AXLE}.

\bibitem[Klypin et al.(2015)]{Klypin2015}
Klypin, A., Karachentsev, I., Makarov, D., \& Nasonova, O. (2015). \emph{Abundance of field galaxies}. MNRAS, 454, 1798--1810.

\bibitem[Lotz et al.(2017)]{Lotz2017}
Lotz, J. M., et al.\ (2017). \emph{The Frontier Fields: Survey design and initial results}. ApJ, 837, 97.

\bibitem[Natarajan, Chiang \& Dutra(2026)]{NatarajanChiangDutra2026}
Natarajan, P., Chiang, B. T., \& Dutra, I. (2026). \emph{Stress-testing the cold dark matter paradigm in galaxy clusters with strong lensing}. The Astrophysical Journal Letters. Published 14 May 2026. doi:\href{https://doi.org/10.3847/2041-8213/ae53ea}{10.3847/2041-8213/ae53ea}.

\bibitem[Navarro, Frenk \& White(1997)]{NFW1997}
Navarro, J. F., Frenk, C. S., \& White, S. D. M. (1997). \emph{A universal density profile from hierarchical clustering}. ApJ, 490, 493--508.

\bibitem[Tokayer et al.(2024)]{Tokayer2024}
Tokayer, Y. M., et al.\ (2024). \emph{Baryonic feedback in galaxy cluster simulations: limits on the inner profile}. ApJ, 967, 132.

\bibitem[Yale News(2026)]{YaleNews2026}
Shelton, J. (2026). \emph{Is it time to expand our thinking about dark matter? A new study says yes}. Yale News, 14 May 2026. \url{https://news.yale.edu/2026/05/14/it-time-expand-our-thinking-about-dark-matter-new-study-says-yes}.

\end{thebibliography}

\vspace{1em}
\noindent\rule{\textwidth}{0.4pt}
\noindent\footnotesize\textit{Manuscript history.} V1 submitted 2026-06-27 to \emph{Open Journal of Astrophysics} (Scholastica manuscript 3187717). V2 (this document) corrects the central observational citation (Natarajan, Chiang \& Dutra 2026, ApJL, doi:10.3847/2041-8213/ae53ea, full author list and journal), expands the comparison to the Yale group's two proposed resolutions, and adds the AXLE Lean~4 verification anchor. The dm\textsuperscript{3} framework's mathematical content and numerical predictions are unchanged.

\end{document}
